\exercisesection{\S~1}

\begin{enumerate}
\def\labelenumi{\arabic{enumi})}
\item
  Assume that \(k\) has characteristic \(p > 0\) . Let \(V\) be a vector space, \(S\) a finite set. A map \(r: S \to \operatorname{End}(V)\) satisfies condition (AC) if and only if there exists a power \(q\) of \(p\) such that \([s^q, s'^q] = 0\) for all \(s, s' \in S\) . (Use Chap. I, §1, Exerc. 19, formula (1).)
\item
  Assume that \(k\) is perfect. Let \(V\) be a finite dimensional vector space, and \(u, v \in \operatorname{End}(V)\) . Let \(u_s, u_n, v_s, v_n\) be the semi-simple and nilpotent components of \(u, v\) . The following conditions are equivalent: (i) there exists an integer \(m\) such that \((\operatorname{ad} u)^m v = 0\) ; (ii) \(u_s\) and \(v\) commute. (To prove (i) \(\Longrightarrow\) (ii), reduce to the case where \(k\) is algebraically closed and use Lemma 1 (ii).)
\item
  We make the assumptions in no. 2. Assume that \(k\) is infinite and that condition (AC) is satisfied. Let \(k'\) be a perfect extension of \(k\) . Let \(\lambda : S \to k\) be such that \(V^{\lambda}(S) \neq 0\) . Put
\end{enumerate}

\[
\mathrm{V} ^ {\prime} = \mathrm{V} \otimes_ {k} k ^ {\prime}, \quad \mathrm{S} ^ {\prime} = \mathrm{S} \otimes_ {k} k ^ {\prime}.
\]

Let \(r' : S' \to \text{End}(V')\) be the linear map obtained from \(r\) by extension of scalars. There exists a unique map \(\lambda' : S' \to k'\) such that \(V^{\lambda}(S) \otimes_k k' = V'^{\lambda'}(S')\) . (Reduce to the case where \(V = V^{\lambda}(S)\) . Let \(P\) be a polynomial function on \(S\) and \(q\) a power of the characteristic exponent of \(k\) dividing \(\dim V\) , such that \(\lambda^q = P\) . Let \(P'\) be the polynomial function on \(S'\) which extends \(P\) . For each \(s' \in S'\) , there exists a \(\lambda'(s') \in k'\) such that \(\lambda'(s')^q = P'(s')\) . Show that the characteristic polynomial of \(r'(s')\) is \((X - \lambda'(s'))^{\dim V}\) .)

\begin{enumerate}
\def\labelenumi{\arabic{enumi})}
\setcounter{enumi}{3}
\tightlist
\item
  Assume that k has characteristic zero. Let \(\mathfrak{g} = \mathfrak{sl}(3, k)\) and let a be the subalgebra of g generated by a diagonal matrix with eigenvalues 1, -1, 0. Show that a is reductive in g, that the commutant m of a in g consists of the diagonal matrices of trace zero, and that the commutant of m in g is equal to m, and hence is distinct from a (cf.~no. 5, Remark).
\end{enumerate}

¶ 5) Assume that k is infinite. Let g be a Lie algebra and V a g-module. If n is an integer \(\geq 0\) , denote by \(V_{n}\) the set of \(v \in V\) such that \(x^{n}v = 0\) for all \(x \in g\) .

\begin{enumerate}
\def\labelenumi{\alph{enumi})}
\tightlist
\item
  Show that, if \(v \in V_n\) , \(x, y \in \mathfrak{g}\) , then
\end{enumerate}

\[
\left(\sum_ {i = 1} ^ {n} x ^ {n - i} y x ^ {i - 1}\right) v = 0.
\]

(Use the fact that \((x + ty)^n v = 0\) for all \(t\in k\) .)

Replacing \(y\) by \([x, y]\) in this formula, deduce \(^3\) that \((x^n y - y x^n)v = 0\) , and hence that \(x^n y v = 0\) .

\begin{enumerate}
\def\labelenumi{\alph{enumi})}
\setcounter{enumi}{1}
\item
  Show that \(V_{n}\) is a g-submodule of V (use a)). In particular \(\mathrm{V}^{0}(\mathfrak{g})=\bigcup_{n}\mathrm{V}_{n}\) is a g-submodule of V.
\item
  Assume that \(k = \mathbf{R}\) or \(\mathbf{C}\) , and denote by \(G\) a simply-connected Lie group with Lie algebra \(\mathfrak{g}\) ; the action of \(\mathfrak{g}\) on \(V\) defines a law of operation of \(G\) on \(V\) (Chap. III, §6, no. 1). Show that an element \(v \in V\) belongs to \(V_n\) if and only if \((s - 1)^n v = 0\) for all \(s \in G\) ; in particular \(V^0(\mathfrak{g}) = V^1(G)\) .
\end{enumerate}

\begin{enumerate}
\def\labelenumi{\arabic{enumi})}
\setcounter{enumi}{5}
\tightlist
\item
  The notations are those of Exerc. 12 of Chap. I, §3. In particular, g is a Lie algebra, M a g-module, and \(\mathrm{H}^{p}(\mathfrak{g},\mathrm{M})=\mathrm{Z}^{p}(\mathfrak{g},\mathrm{M})/\mathrm{B}^{p}(\mathfrak{g},\mathrm{M})\) is the cohomology space of degree p of g with values in M.
\end{enumerate}

\begin{enumerate}
\def\labelenumi{\alph{enumi})}
\item
  Show that \(\mathrm{B}^{p}(\mathfrak{g},\mathrm{M})\) and \(\mathrm{Z}^{p}(\mathfrak{g},\mathrm{M})\) are stable under the natural representation \(\theta\) of g on the space of cochains \(\mathrm{C}^{p}(\mathfrak{g},\mathrm{M})\) . It follows that there is a representation of g on \(\mathrm{H}^{p}(\mathfrak{g},\mathrm{M})\) . Show that this representation is trivial (use the formula \(\theta = di + id\) , loc. cit.).
\item
  Let \(\overline{k}\) be an algebraic closure of k. Let \(x \in g\) and let \(x_{M}\) be the corresponding endomorphism of M. Let \(\lambda_{1}, \ldots, \lambda_{n}\) (resp. \(\mu_{1}, \ldots, \mu_{m}\) ) be the eigenvalues (in \(\overline{k}\) ) of \(ad_{g}x\) (resp. of \(x_{M}\) ), repeated according to their multiplicity. Show that the eigenvalues of the endomorphism \(\theta(x)\) of \(C^{p}(g, M)\) are the \(\mu_{j} - (\lambda_{i_{1}} + \cdots + \lambda_{i_{p}})\) , where \(1 \leq j \leq m\) and
\end{enumerate}

\[
1 \leq i _ {1} <   i _ {2} <   \dots <   i _ {p} \leq n.
\]

Deduce, using a), that \(\mathrm{H}^{p}(\mathfrak{g},\mathrm{M})=0\) if none of the \(\mu_{j}-(\lambda_{i_{1}}+\cdots+\lambda_{i_{p}})\) is zero.

\begin{enumerate}
\def\labelenumi{\alph{enumi})}
\setcounter{enumi}{2}
\tightlist
\item
  Assume that the representation \(\mathfrak{g} \to \operatorname{End}(\mathrm{M})\) is faithful, and that \(x\) satisfies the condition:
\end{enumerate}

\[
\mu_ {j _ {1}} + \dots + \mu_ {j _ {p}} \neq \mu_ {k _ {1}} + \dots + \mu_ {k _ {p + 1}}\tag{($S_{p}$)}
\]

for all \(j_{1},\ldots,j_{p},k_{1},\ldots,k_{p+1}\in[1,m]\) .

Show that we then have \(\mathrm{H}^p (\mathfrak{g},\mathrm{M}) = 0\) (remark that the eigenvalues \(\lambda_{i}\) of \(\mathrm{ad}_{\mathfrak{g}}x\) are of the form \(\mu_j - \mu_k\) , and apply \(b\) )).

\begin{enumerate}
\def\labelenumi{\arabic{enumi})}
\setcounter{enumi}{6}
\tightlist
\item
  Let \(\mathfrak{g}\) be a nilpotent Lie algebra and V a \(\mathfrak{g}\) -module such that \(V^0(\mathfrak{g}) = 0\) . Show that \(H^p(\mathfrak{g}, V) = 0\) for all \(p \geq 0\) . (Reduce to the case where \(V = V^\lambda(\mathfrak{g})\) , with \(\lambda \neq 0\) and choose an element \(x \in \mathfrak{g}\) such that \(\lambda(x) \neq 0\) . Apply Exerc. 6 b), remarking that the \(\lambda_i\) are all zero and the \(\mu_j\) are all equal to \(\lambda(x)\) .)
\end{enumerate}

Recover the Cor. of Prop. 9 (take p = 1). \(^{4}\)

¶ 8) Assume that k is of characteristic p \textgreater{} 0. Let g be a Lie algebra over k with basis \((e_{1},\ldots,e_{n})\) . Denote by U the enveloping algebra of g and C the centre of U. For \(i = 1,\ldots,n\) , choose a non-zero p-polynomial \(f_{i}\) , of degree \(d_{i}\) , such that \(f_{i}(\mathrm{ad}e_{i}) = 0\) ; then \(f_{i}(e_{i}) \in \mathrm{C}\) , cf.~Chap. I, §7, Exerc. 5. Put \(z_{i} = f_{i}(e_{i})\) .

\begin{enumerate}
\def\labelenumi{\alph{enumi})}
\item
  Show that \(z_{1},\ldots,z_{n}\) are algebraically independent. If \(A=k[z_{1},\ldots,z_{n}]\) , show that U is a free A-module with basis the monomials \(e_{1}^{\alpha_{1}}\ldots e_{n}^{\alpha_{n}}\) , where \(0\leq\alpha_{i}\leq d_{i}\) . (Use the Poincaré-Birkhoff-Witt theorem.) The rank {[}U:A{]} of U over A is equal to \(d_{1}\ldots d_{n}\) ; it is a power of p.~Deduce that C is an A-module of finite type, hence a k-algebra of finite type and of dimension n (Commutative Algebra, Chap. VIII).
\item
  Let \(\mathrm{K}\) be the field of fractions of \(\mathrm{A}\) , and let
\end{enumerate}

\[
\mathrm{U} _ {(\mathrm{K})} = \mathrm{U} \otimes_ {\mathrm{A}} \mathrm{K}, \quad \mathrm{C} _ {(\mathrm{K})} = \mathrm{C} \otimes_ {\mathrm{A}} \mathrm{K}.
\]

Then \(\mathrm{U}_{(\mathrm{K})}\supset\mathrm{C}_{(\mathrm{K})}\supset\mathrm{K}\) . Show that \(\mathrm{U}_{(\mathrm{K})}\) is a field with centre \(\mathrm{C}_{(\mathrm{K})}\) , and that this is the quotient field (both left and right) of U, cf.~Chap. I, §2, Exerc. 10. Deduce that \([\mathrm{U}_{(\mathrm{K})}:\mathrm{C}_{(\mathrm{K})}]\) is of the from \(q^{2}\) , where q is a power of p; we have \([\mathrm{C}_{(\mathrm{K})}:\mathrm{K}]=q_{\mathrm{C}}\) , where \(q_{C}\) is a power of p, and \([U:A]=q_{C}q^{2}\) .

\begin{enumerate}
\def\labelenumi{\alph{enumi})}
\setcounter{enumi}{2}
\tightlist
\item
  Let d be a non-zero element of A, and let \(\Lambda\) be a subring of \(\mathrm{U}_{(\mathrm{K})}\) such that \(U \subset \Lambda \subset d^{-1}U\) . Show that \(\Lambda = U\) . {[}If x = b/a, \(a \in A - \{0\}\) , is an element of \(\Lambda\) , show by induction on m that the relation \(b \in Ua + U_m\) implies that \(b \in Ua + U_{m-1}\) , where \(\{U_m\}\) is the canonical filtration of U. (For this, use the fact that gr U is integrally closed, and argue as in Prop. 15 of Commutative Algebra, Chap. V, §1, no. 4.) For m = 0, this gives \(b \in Ua\) , i.e.~\(x \in U\) .{]}
\end{enumerate}

Deduce that C is integrally closed.

\begin{enumerate}
\def\labelenumi{\alph{enumi})}
\setcounter{enumi}{3}
\tightlist
\item
  Assume that k is algebraically closed. Let \(\rho: \mathfrak{g} \to \mathfrak{gl}(V)\) be an irreducible linear representation of g and \(\rho_{U}\) the corresponding representation of U. The restriction of \(\rho_{U}\) to C is a homomorphism \(\gamma_{\rho}\) from C to k (identified with the homotheties of V); let \(\alpha_{\rho}\) be its restriction to A. Show that for any homomorphism \(\alpha\) (resp. \(\gamma\) ) from the k-algebra A (resp. C) to k, there exists at least one irreducible representation \(\rho\) of g such that \(\alpha_{\rho} = \alpha\) (resp. \(\gamma_{\rho} = \gamma\) ) and that there are only finitely-many such representations (up to equivalence). Show that \(\dim V \leq q\) , with the notation of b). \(^{5}\)
\end{enumerate}

¶ 9) We retain the notations of the preceding exercise, and assume further that g is nilpotent.

\begin{enumerate}
\def\labelenumi{\alph{enumi})}
\item
  Show that the basis \((e_{1},\ldots,e_{n})\) can be chosen so that, for any pair \((i,j)\) , \([e_{i},e_{j}]\) is a linear combination of the \(e_{h}\) for \(h > \sup(i,j)\) . Assume from now on that the \(e_{i}\) satisfy this condition. For \(i=1,\ldots,n\) choose a power \(q(i)\) of p such that \(\operatorname{ad}(e_{i})^{q(i)} = 0\) , and put \(z_{i} = e_{i}^{q(i)}, A = k[z_{1},\ldots,z_{n}]\) , cf.~Exerc. 8. b) Let \(\rho : g \to gl(V)\) be a linear representation of g. Assume that \(\rho(e_{i})\) is nilpotent for \(i = 1,\ldots,n\) . Show that \(\rho(x)\) is nilpotent for all \(x \in g\) . (Argue by induction on \(n = \dim g\) and reduce to the case where \(\rho\) is irreducible. Show that, in this case, \(\rho(e_{n}) = 0\) and apply the induction hypothesis.)
\item
  Let \(\rho_1: \mathfrak{g} \to \mathfrak{gl}(V_1)\) and \(\rho_2: \mathfrak{g} \to \mathfrak{gl}(V_2)\) be two linear representations of \(\mathfrak{g}\) . Assume that \(V_1\) and \(V_2\) are \(\neq 0\) , and that \(V_1 = V^{\lambda_1}(\mathfrak{g}), V_2 = V^{\lambda_2}(\mathfrak{g})\) , where \(\lambda_1\) and \(\lambda_2\) are two functions on \(\mathfrak{g}\) , cf.~no. 3. Show that, if \(\lambda_1(e_i) = \lambda_2(e_i)\) for \(i = 1, \ldots, n\) , then \(\lambda_1 = \lambda_2\) and there exists a non-zero \(\mathfrak{g}\) -homomorphism from \(V_1\) to \(V_2\) (apply \(b\) ) to the \(\mathfrak{g}\) -module \(V = \mathcal{L}(V_1, V_2)\) and use Engel's theorem to show that \(V\) contains a non-zero \(\mathfrak{g}\) -invariant element). Deduce that if, in addition, \(V_1\) and \(V_2\) are simple, they are isomorphic.
\item
  Assume that k is algebraically closed. Let R be the set of equivalence classes of irreducible representations of g. If \(\rho \in R\) , put
\end{enumerate}

\[
x _ {\rho} = (x _ {\rho} (1), \dots , x _ {\rho} (n)) \in k ^ {n},
\]

where \(x_{\rho}(i)\) is the unique eigenvalue of \(\rho(e_{i})\) . Show that \(\rho \mapsto x_{\rho}\) is a bijection from R to \(k^{n}\) . (Injectivity follows from c), and surjectivity from Exerc. 8 d). Deduce the following consequences:

\begin{enumerate}
\def\labelenumi{(\roman{enumi})}
\item
  For any maximal ideal \(\mathfrak{m}\) of A, the quotient of U/mU by its radical is a matrix algebra.
\item
  The degree of any irreducible representation of \(\mathfrak{g}\) is a power of \(p\) (this follows from (i) and the fact that \([\mathrm{U} / \mathfrak{mU}:k]\) is a power of \(p\) ).
\item
  Every homomorphism from A to k extends uniquely to a homomorphism from C to k (use the fact that C/mC is contained in the centre of U/mU, which is a local k-algebra with residue field k).
\item
  There exists an integer N ≥ 0 such that \(x^{p^{N}} \in A\) for all \(x \in C\) (this follows from (iii)). \(^{6}\)
\end{enumerate}

¶ 10) Assume that \(k\) is of characteristic \(p > 0\) . Denote by \(\mathfrak{g}\) a Lie algebra with basis \(\{e_1, e_2, e_3\}\) , with \([e_1, e_2] = e_3\) , \([e_1, e_3] = [e_2, e_3] = 0\) .

\begin{enumerate}
\def\labelenumi{\alph{enumi})}
\item
  Show that the centre of \(\mathrm{U}\mathfrak{g}\) is \(k[e_1^p,e_2^p,e_3]\) .
\item
  Assume that k is algebraically closed. Show that, for all \((\lambda_{1},\lambda_{2},\lambda_{3})\in k^{3}\) , there exists (up to equivalence) a unique irreducible representation \(\rho\) of g such that \(\lambda_{i}\) is the unique eigenvalue of \(\rho(e_i)\) ( \(i = 1,2,3\) ); the degree of \(\rho\) is \(p\) if \(\lambda_3 \neq 0\) , and is 1 if \(\lambda_3 = 0\) . (Apply Exerc. 8 and 9, or argue directly.)
\end{enumerate}

\begin{enumerate}
\def\labelenumi{\arabic{enumi})}
\setcounter{enumi}{10}
\tightlist
\item
  Let \(\mathfrak{h}\) be a nilpotent Lie algebra, \(V\) an \(\mathfrak{h}\) -module not reduced to 0, and \(\lambda\) a function on \(\mathfrak{h}\) such that \(V = V^{\lambda}(\mathfrak{h})\) . Prove the equivalence of the following properties:
\end{enumerate}

\begin{enumerate}
\def\labelenumi{(\roman{enumi})}
\item
  \(\lambda\) is a linear form on \(\mathfrak{h}\) , zero on \([\mathfrak{h},\mathfrak{h}]\) .
\item
  There exists a basis of \(\mathrm{V}\) with respect to which the endomorphisms defined by the elements of \(\mathfrak{h}\) are triangular.
\end{enumerate}

(To prove that (i) \(\Longrightarrow\) (ii), apply Engel's theorem to the \(\mathfrak{h}\) -module \(\mathcal{L}(\mathrm{W}_{\lambda},\mathrm{V})\) , where \(\mathrm{W}\) is the 1-dimensional \(\mathfrak{h}\) -module defined by \(\lambda\) .)

Properties (i) and (ii) are true if k is of characteristic 0 (Prop. 9).

